\parttitle{Mathematical destination}
\labelled{Reversing the finish needs care.}{In last-counter-loses Nim, start with at least one counter in total; individual piles may be empty. A move removes any positive number from one nonempty pile. If all nonempty piles have size one, the next player loses exactly when their number is odd. If at least one pile exceeds one, the losing positions are the ordinary zero-xor positions. An ordinary Nim strategy must change at the move leaving only singletons: leave an odd number. An empty pile beside a nonempty pile is a legal start, distinct from the all-empty position: the latter is not played, because the player who just emptied the table has lost.}
\labelled{Coins turn gaps into piles.}{In the no-jumping coin row, number the playable squares 1,2,$\ldots$; the left wall is before square 1. Pair coins from the right. Record empty squares inside each pair, plus the gap before an unpaired leftmost coin if there is one. The position loses exactly when these gap sizes have xor zero. For two coins, adjacency alone decides a losing start. For three coins at $x_1<x_2<x_3$, the relevant gaps are $x_1-1$ and $x_3-x_2-1$.}
\labelled{Two games can cancel.}{On each turn choose one rook board and slide its token left or down. The player who leaves both tokens at their stars wins. The four distances to the stars act as four Nim piles. Two identical starts lose by copying on the other board. Distinct starts also lose when their two coordinate-xor values agree: individually winning components can combine to a losing game. Adding two independently losing components always remains losing.}
\labelled{What children discover.}{The end rule, physical order and choice of component matter. Repeated wins are experiments; a guarantee needs a reply strategy or invariant. Binary bundles from base Week 8 are an available explanation, not a prerequisite for playing. These investigations extend rather than repeat its ordinary three-pile enumeration and queen game.}
\labelled{Band map.}{Young children can slide two coins on a long strip and test adjacency; a partner enforces no jumping. The printed coin entry is K--3, its three-coin continuation 2--5, last-loses page 2--5, and two-board pages 4--5. Younger ready players can join acted two-board games with an adult. Grades 4--5 can explain the finishing exception, paired gaps and cancellation. Counting to 12, subtraction and recognising neighbours suffice for entry. Accept an acted reply rule before symbolic xor.}
\labelled{Status.}{Prepared, unpiloted, with no physical material-fit rehearsal. A first visit can stay with one investigation for 20--30 minutes or more. The guide retains deeper work and later visits; finishing the companion is not the goal.}
\newpage
\parttitle{Materials, launch and first route}
\labelled{Per pair.}{About 24 counters at most 0.65 inch across, three pile areas, a fixed 12-square coin strip and two distinguishable rook tokens at most 0.75 inch across. The largest printed pile has four counters; four 0.65-inch counters can lie in a row inside a 2.9-by-1.4-inch pile mat. Larger invented piles may use the tabletop. The five-page student companion is landscape Letter. Coin working strips on pages 2--3 have twelve 0.85-inch squares (10.2 inches). Page 4 has two independent 4-by-4 rook boards with 0.85-inch squares (3.4 inches square), one token each. Smaller coin and rook diagrams are start references, not active playing boards. Print at 100\% and choose pieces that fit comfortably within each square. A table strip with 0.85-inch squares is 10.2 inches long; it may be made with tape rather than shrunk to one Letter page. All counts and fit suggestions are preparation estimates, not a physical pretest.}
\labelled{Visible conventions.}{Keep the coin strip's left end fixed and label its first playable square 1. A coin cannot jump another or leave the strip. In the two-board game a star token stays put; the other board can still be played. Finish only when both tokens are at their stars. In misere play announce the changed final-counter rule before starting a new game, and keep a rule card on the table.}
\labelled{Launch together, 3--5 minutes.}{Allow free handling, then show one legal coin slide into an empty square and one attempted jump that the partner rejects. In a two-board practice game move only one token and let a child make the next legal move on either board. Demonstrate the different ending with a one-counter pile before misere play. Introduce each convention when choosing its investigation; avoid combining all rules in one oral introduction.}
\labelled{A manageable route.}{For a concrete first visit use coin slides, taking 15--25 minutes for several starts and partner-created starts. A later visit may explore the end-rule change, then two-board cancellation if children retain each component's rules. The current three adults remain anchored at KK11 / 3333 / 445; the upper trio can rotate players and referee. Adults may read or scribe, while children choose the starts and moves.}
\labelled{Hints and records.}{Leave actual coin positions visible. One start and one replayable move sequence is enough. Later, a gap row can help compare different layouts, but an adult should not impose it before children have noticed what changes. Record the exact end rule and which starts were tried. Conjectures about all positions should be marked separately from checked finite instances.}
\labelled{Sources and limits.}{These are original task constructions using standard Nim and paired-gap arguments. Current Week 8 F08-v4 provides normal two/three-pile Nim and binary bundles; current Week 7 already has misere subtraction, a related but different game. \emph{Math Circle by the Bay}, preface printed pp. viii--x/PDF 9--11, reports manipulatives, varying depth and reserve challenges; our arrangements and timing are adaptations. No classroom observations are available.}
\newpage
\parttitle{Problem 1 / Taking the last counter loses}
\labelled{Printed starts.}{On page 1, starts 1--12 favour respectively second, first, first, second, first, second, first, second, first, first, second, first. Winning first moves in starts 2,3,5,7,9,10,12 can leave respectively (1), (1), (2,2), (1), (1,1,1), (2,2), (2,3,1). Start 7 is (2,0), a first-player win by leaving one, contrasting with start 1 (1,0). The three-pile contrasts include (1,1,1), which loses under this ending, and (1,1,2), which wins by reducing the large pile to one, leaving three singletons for the opponent.}
\labelled{Small boundaries.}{One pile of one loses. One pile larger than one wins by leaving exactly one. Two piles (1,1) win: take either singleton and leave one for the opponent. Equal piles (2,2), (3,3),$\ldots$ lose. Unequal two-pile starts with a larger pile win: if the small pile is at least two, equalise; if it is one, remove the larger pile completely, leaving one. Empty piles are omitted.}
\labelled{General rule.}{When every surviving pile is a singleton, each turn removes exactly one, so parity decides: odd counts lose, even counts win. When at least two piles exceed one, a move to ordinary xor zero still has a large pile, and the usual Nim invariant applies. If only one large pile remains, its size determines which final singletons to leave: reduce it to one when the existing singleton count is even, or remove it when that count is odd. In both cases the opponent receives an odd singleton count.}
\labelled{Why this is complete.}{A zero-xor position with a large pile cannot have exactly one large pile: the xor of the singleton piles is only 0 or 1, which cannot equal a size at least two. Thus every move from a zero-xor position either leaves a nonzero-xor position with a large pile, or crosses to a winning even-singleton state. Every nonzero-xor position with at least two large piles can make the ordinary zero-xor move; the one-large-pile case uses the boundary strategy above. No losing state moves to a losing state, and every winning state has a move to one.}
\labelled{Entry and hint.}{Play small starts with a visible ``last loses'' card. Ask what happens when all remaining piles contain just one. Let children find why blindly copying until the table is empty fails. If binary reasoning is unfamiliar, the complete two-pile strategy above supplies a substantial concrete investigation by itself.}
\labelled{Continuation.}{Compare a familiar normal-play start with exactly the same counters under this ending. Seek starts whose preferred first/second choice changes and starts whose choice remains. A child's reply strategy can use physical matching and a final singleton adjustment; formal binary notation is optional. Save both the start and the ending rule.}
\newpage
\parttitle{Problem 2 / Coins cannot jump}
\labelled{Printed starts and conversion.}{Page 2 starts 1--6 favour second, first, second, first, second, first. Page 3 starts 7--14 favour second, first, second, first, second, first, second, first. Their relevant gap pairs $(a,b)$ are respectively $(0,0)$, $(0,1)$, $(1,1)$, $(1,2)$, $(2,2)$, $(2,3)$, $(3,3)$, $(3,4)$. The non-task visual moves square 4 to 2 while the coin at 10 stays fixed; no jump occurs.}
\labelled{Two coins.}{If they are adjacent, the second player copies every slide of the left coin by sliding the right coin the same distance, keeping adjacency. Eventually the left coin is against the wall and neither can move. If they have a gap, the first player closes it by sliding the right coin immediately beside the left, handing over an adjacent pair. The unrecorded space to the wall does not decide a two-coin position.}
\labelled{Three coins.}{Let $a=x_1-1$ be the empty squares before the first coin, and $b=x_3-x_2-1$ the empty squares between the right pair. A start loses exactly when $a=b$. If unequal, reduce the larger recorded gap to match the smaller: slide the first coin left to reduce $a$, or the rightmost coin left to reduce $b$. Sliding the middle coin increases $b$; from equality, reply by sliding the rightmost coin left the same distance. The coins remain ordered and no reply jumps one.}
\labelled{More coins and the invariant.}{Pair from the right. Moving a right member decreases one recorded gap; moving a left member increases that pair's gap. An unpaired far-left coin decreases its wall gap. Thus every move changes exactly one recorded value, taking a zero-xor list to nonzero xor. From nonzero xor an ordinary Nim reduction of a chosen gap is always physically legal: move its right coin, or the unpaired leftmost coin, by the reduction. This reaches zero xor. The gaps between pairs are deliberately unrecorded.}
\labelled{Why the drawing bridge matters.}{A coin layout is not an unrestricted pile game until the pairing is explained. Put small arcs below actual paired coins, then place one empty-square marker per counted gap. Only after this action introduce a gap number row if useful. Do not count the distance from the wall to the first coin when the total number of coins is even.}
\labelled{Hint and continuation.}{Ask children to compare adjacent starts far from and near the wall. For three coins ask which two gaps could be made equal. Older children can find several very different losing layouts, or move a left member of a pair to challenge the reply rule. Younger children can remain with two coins and invent starts for a partner. Save the actual square positions, not just the sum of them.}
\newpage
\parttitle{Problem 3 / Two races, one turn}
\labelled{Printed pairs.}{Starts 1--6 favour second, second, second, first, second, first. In coordinates measured right/up from each star they are $(1,2;1,2)$, $(1,1;2,2)$, $(1,0;2,3)$, $(1,3;2,2)$, $(0,2;1,3)$, $(0,1;1,3)$. The worked non-task turns change A from $(2,2)$ to $(0,2)$, then B from $(1,3)$ to $(1,1)$; only one board moves each turn. A further distinct cancelling pair is $(0,3;1,2)$, both individually first-player wins with value 3.}
\labelled{The state.}{For tokens at distances $(x,y)$ and $(u,v)$ from their stars, the position loses exactly when $x\mathbin{\mathrm{xor}}y\mathbin{\mathrm{xor}}u\mathbin{\mathrm{xor}}v=0$. Moving one token changes exactly one coordinate by reducing it, just like one of four piles. A finished token contributes two zero coordinates and stays still.}
\labelled{A concrete cancellation.}{Two identical off-diagonal starts separately favour the first player but combine to favour the second: copy each move on the other board. The board chosen for the reply is different from the board of the opponent's move. Distinct starts can also cancel: $(1,0)$ and $(2,3)$ both have coordinate xor 1, so their combined xor is zero. The first component is a one-step race; the second has several moves. Equality of drawings is sufficient, not necessary.}
\labelled{Two losing components.}{An ordinary losing rook start lies on its diagonal, so $x=y$ and its coordinate xor is zero. Two such starts together still have xor zero. The proposition that combining two individually losing races could create a first-player win is false under these exact independent rules. If a game couples moves across both boards, this conclusion need not hold.}
\labelled{Same totals do not decide.}{A token at $(0,2)$ has coordinate xor 2, while a token at $(1,1)$ has xor 0; both distances total two. Pair either with $(2,0)$. The first combined position has xor zero and loses, the second has xor 2 and wins. Total distance alone misses how the moves are distributed.}
\labelled{Reply proof without notation.}{Identical boards support direct copying. For distinct starts, children may list all small successor pairs, marking a state losing when none of its legal successors lose, and winning when one does. Binary bundles give the more general proof: at zero xor, a move changes the highest altered bundle parity; from nonzero xor, reduce a coordinate containing the highest odd bundle to make all bundle counts even. The token slide performs exactly that reduction.}
\labelled{Hint and return work.}{Ask whether winning each race separately says who wins when turns are shared. Try matching starts, then search for a distinct cancelling pair rather than more identical copies. Keep the two-board state together in the record. Further boards provide later visits, but three ordinary Nim piles are already base Week 8 and are not counted again here.}
